\section{Sliding window}

Maintain a window \texttt{[left, right]} over an array or string, expanding or shrinking it while preserving a condition. Typical uses include longest or shortest valid subarrays, at most $k$ distinct values, fixed-size maximum/minimum sums, and frequency-based substring problems.

\subsection{Variable-size window}
This example finds the longest subarray with sum at most $k$. It requires all values to be non-negative.

\begin{lstlisting}
int n, k;
cin >> n >> k;
vector<int> a(n);
for (int &x : a) cin >> x;

long long sum = 0;
int left = 0, ans = 0;
for (int right = 0; right < n; right++) {
    sum += a[right];

    // Shrink until the window becomes valid again.
    while (left <= right && sum > k) {
        sum -= a[left];
        left++;
    }

    // [left, right] is valid.
    ans = max(ans, right - left + 1);
}
cout << ans << '\n';
\end{lstlisting}

It is usually $O(n)$ time and uses $O(k)$ or $O(\text{alphabet size})$ memory, depending on the data structure used.

\subsection{Fixed-size window}
Use this for a window containing exactly $k$ elements: maximum/minimum subarray sum, maximum average, or a count/frequency calculation for every window.

\begin{lstlisting}
long long windowSum = 0;
for (int i = 0; i < n; i++) {
    windowSum += a[i];
    if (i >= k) windowSum -= a[i - k];
    if (i >= k - 1) {
        // [i-k+1, i] is the current window. Process it here.
    }
}
\end{lstlisting}
This pattern is $O(n)$.

\section{Sparse table}

Preprocess a static array for fast range queries with an idempotent operation, such as range minimum, maximum, or gcd. For RMQ, query \texttt{[l, r]} with two overlapping blocks. Build time and memory are $O(n \log n)$; each query is $O(1)$. For frequently updated values, use a segment tree instead.

\subsection{Range minimum query}
\begin{lstlisting}
class SparseTable {
private:
    int n, LOG;
    vector<vector<int>> st;
    vector<int> lg;

public:
    SparseTable(const vector<int>& a) {
        n = a.size();
        lg.resize(n + 1);
        for (int i = 2; i <= n; i++) lg[i] = lg[i / 2] + 1;
        LOG = lg[n] + 1;
        st.assign(LOG, vector<int>(n));

        // Base level: intervals of length 1.
        for (int i = 0; i < n; i++) st[0][i] = a[i];

        // st[j][i] is the minimum of a[i ... i + 2^j - 1].
        for (int j = 1; j < LOG; j++)
            for (int i = 0; i + (1 << j) <= n; i++)
                st[j][i] = min(st[j - 1][i],
                               st[j - 1][i + (1 << (j - 1))]);
    }

    // Minimum in the inclusive range [l, r].
    int query(int l, int r) {
        int len = r - l + 1;
        int j = lg[len];
        return min(st[j][l], st[j][r - (1 << j) + 1]);
    }
};

int n, q;
cin >> n >> q;
vector<int> a(n);
for (int &x : a) cin >> x;
SparseTable sp(a);
while (q--) {
    int l, r;
    cin >> l >> r;
    // The example is 0-indexed. For 1-indexed input: --l; --r;
    cout << sp.query(l, r) << '\n';
}
\end{lstlisting}

For range maximum, replace \texttt{min} with \texttt{max}. For range gcd, replace it with \texttt{gcd}; the preprocessing structure is unchanged.

\section{Binary exponentiation}

Compute $a^b$ in $O(\log b)$ using repeated squaring, instead of multiplying $a$ by itself $b$ times. It also applies to matrix exponentiation and any associative operation that supports repeated squaring. The result can overflow if it does not fit in the selected type.

\begin{lstlisting}
long long binpow(long long a, long long b) {
    long long ans = 1;
    while (b > 0) {
        if (b & 1) ans *= a;
        a *= a;
        b >>= 1;
    }
    return ans;
}

long long a, b;
cin >> a >> b;
cout << binpow(a, b) << '\n';
\end{lstlisting}

\section{Modular binary exponentiation}

Compute $a^b \bmod m$ without forming the potentially enormous value $a^b$. It is useful in modular arithmetic, combinatorics, number theory, and modular inverse calculations with a prime modulus: $a^{m-2} \bmod m$. The \texttt{\_\_int128} multiplication prevents overflow before reducing modulo \texttt{mod}.

\begin{lstlisting}
long long binpowmod(long long a, long long b, long long mod) {
    long long ans = 1 % mod;
    a %= mod;
    while (b > 0) {
        if (b & 1) ans = (__int128)ans * a % mod;
        a = (__int128)a * a % mod;
        b >>= 1;
    }
    return ans;
}

long long a, b, mod;
cin >> a >> b >> mod;
cout << binpowmod(a, b, mod) << '\n';
\end{lstlisting}

For input \texttt{2 10 1000}, the result is \texttt{24}, because $2^{10} = 1024$ and $1024 \bmod 1000 = 24$.

\section{Quick reference}
\begin{itemize}
\item \textbf{Sliding window}: contiguous subarray or substring problems; $O(n)$.
\item \textbf{Sparse table}: static range queries; build $O(n \log n)$, query $O(1)$.
\item \textbf{binpow}: fast $a^b$; $O(\log b)$.
\item \textbf{binpowmod}: fast $a^b \bmod m$; $O(\log b)$.
\end{itemize}

\section{Competitive-programming tips}
\begin{itemize}
\item Sliding window works naturally when its condition can be adjusted monotonically. Be careful with negative values.
\item Sparse tables are best for static arrays with many queries.
\item Use \texttt{binpow} when the exponent is large.
\item Prefer \texttt{binpowmod} to calculating $a^b$ first when working under a modulus.
\item The sparse-table example uses zero-indexed inclusive ranges \texttt{[l, r]}.
\end{itemize}
